\section*{Hoofdstuk \ref{ch:b2:funcseq} --- Rijen en reeksen van
functies}

\begin{solution}{exo:b2:funcseq:1}
$f_n(x) = \frac{x}{1 + nx}$: \omterm{def:b2:funcseq:def}{puntsgewijze} limiet $0$ op
$\intcc{0}{1}$. \omterm{def:b2:funcseq:def}{Uniform}: $f_n$ stijgt op $\intcc{0}{1}$ (afgeleide
$\frac{1}{(1+nx)^2} > 0$), dus $\norm{f_n}_\infty = f_n(1) =
\frac{1}{1+n} \to 0$: \emph{\omterm{def:b2:funcseq:def}{uniform}} op $\intcc{0}{1}$.

$g_n(x) = nx\,\eu^{-nx^2}$: \omterm{def:b2:funcseq:def}{puntsgewijze} limiet $0$ (de
exponentiële wint). Supremum: $g_n' = n\eu^{-nx^2}(1 - 2nx^2)$
verdwijnt in $x_n = \frac{1}{\sqrt{2n}}$, waar $g_n(x_n) =
\sqrt{\frac n2}\,\eu^{-1/2} \to \infty$: dus niet \omterm{def:b2:funcseq:def}{uniform} op
$\intcc{0}{1}$ --- wel \omterm{def:b2:funcseq:def}{uniform} op $\intco{\delta}{1}$, want daar is
$g_n(x) \leq n\,\eu^{-n\delta^2} \to 0$.

$h_n(x) = x^n(1 - x^n)$: \omterm{def:b2:funcseq:def}{puntsgewijze} limiet $0$ op $\intcc{0}{1}$
(beide factoren; in $x = 1$ is $h_n = 0$). Supremum: met $u = x^n
\in \intcc{0}{1}$ is $u(1-u) \leq \frac14$, aangenomen in $u =
\frac12$, dat wil zeggen $x = 2^{-1/n} \in \intoo{0}{1}$:
$\norm{h_n}_\infty = \frac14 \not\to 0$, dus niet \omterm{def:b2:funcseq:def}{uniform} op
$\intcc{0}{1}$; wel \omterm{def:b2:funcseq:def}{uniform} op $\intcc{0}{a}$ ($\sup \leq a^n \to
0$).
\end{solution}

\begin{solution}{exo:b2:funcseq:2}
Er geldt $\int_0^1 nx\,\eu^{-nx^2}\dd x = \bigl[-\tfrac12
\eu^{-nx^2}\bigr]_0^1 = \frac{1 - \eu^{-n}}{2} \to \frac12$,
terwijl $\int_0^1 \lim g_n = 0$. Geen tegenspraak:
\cref{thm:b2:funcseq:integration} vraagt \emph{\omterm{def:b2:funcseq:def}{uniforme}}
convergentie op het segment, en die faalt hier (de bult van hoogte
$\sim\sqrt n$ schuift naar $0$).
\end{solution}

\begin{solution}{exo:b2:funcseq:3}
\omterm{def:b2:funcseq:series}{Normale convergentie} op $\intcc{-1}{1}$: $\norm{x^n/n^2}_\infty =
\frac{1}{n^2}$, \omterm{def:b2:series:summable}{sommeerbaar}: dus is $S$ daar \omterm{def:b2:metric:continuity}{continu}
(\cref{thm:b2:funcseq:seriestransfer}).

Afgeleide: de afgeleide reeks $\sum \frac{x^{n-1}}{n}$ convergeert
\omterm{def:b2:funcseq:series}{normaal} op elke $\intcc{-a}{a}$ met $a < 1$ ($\sup =
\frac{a^{n-1}}{n}$): dus is $S$ van klasse $C^1$ op $\intoo{-1}{1}$
met
\[
S'(x) = \sum_{n\geq1} \frac{x^{n-1}}{n} = \frac1x \sum_{n\geq1}
\frac{x^n}{n} = -\frac{\ln(1 - x)}{x}
\qquad (0 < \abs x < 1),
\]
waarbij de laatste identiteit de logaritmische reeks uit
bachelorjaar 1 is (in \cref{ch:b2:powerseries} eerlijk opnieuw
afgeleid).
\end{solution}

\begin{solution}{exo:b2:funcseq:4}
Bij vaste $x > 0$ is de reeks alternerend met $\frac{1}{n + x}
\downarrow 0$: dus \omterm{def:b2:funcseq:def}{puntsgewijze convergentie}, en de grens voor de
restsom $\abs{R_N(x)} \leq \frac{1}{N + 1 + x} \leq \frac{1}{N+1}$
is \emph{\omterm{def:b2:funcseq:def}{uniform}} op $\intco{\delta}{\infty}$ (zelfs op
$\intoo{0}{\infty}$): \omterm{def:b2:funcseq:def}{uniforme convergentie}. (Niet \omterm{def:b2:funcseq:series}{normaal}:
$\norm{\frac{(-1)^n}{n+x}}_\infty = \frac{1}{n + \delta}$,
divergent.) De \omterm{def:b2:metric:continuity}{continuïteit} volgt uit
\cref{thm:b2:funcseq:seriestransfer}.

Functionaalvergelijking: herindexeer $F(x + 1)$ met $m = n + 1$:
\[
F(x+1) = \sum_{n\geq0} \frac{(-1)^n}{n + 1 + x}
= \sum_{m\geq1}\frac{(-1)^{m-1}}{m+x} ,
\]
dus, na de term $m = 0$ van $F(x)$ af te zonderen,
\[
F(x) + F(x+1)
= \frac{1}{x} + \sum_{m\geq1}
\frac{(-1)^m + (-1)^{m-1}}{m+x} = \frac1x .
\]
\end{solution}

\begin{solution}{exo:b2:funcseq:5}
Zet $g_n = f_n - f \geq 0$ (per hypothese dalend in $n$; de limiet
is puntsgewijs $0$); elke $g_n$ is \omterm{def:b2:metric:continuity}{continu}. Leg $\varepsilon > 0$
vast en zij $U_n = \{x : g_n(x) < \varepsilon\}$: \omterm{def:b2:metric:topology}{open} (origineel
van een \omterm{def:b2:metric:topology}{open} verzameling), stijgend ($g_{n+1} \leq g_n$) en $K$
overdekkend (\omterm{def:b2:funcseq:def}{puntsgewijze convergentie}). Volgens
Borel--Lebesgue (\cref{thm:b2:metric:borellebesgue}) overdekken
eindig veel $U_{n_1} \subseteq \dots \subseteq U_{n_k}$ de ruimte
$K$: dus $K = U_{n_k}$, dat wil zeggen $\norm{g_{n_k}}_\infty \leq
\varepsilon$, en wegens de monotonie $\norm{g_n}_\infty \leq
\varepsilon$ voor alle $n \geq n_k$: \omterm{def:b2:funcseq:def}{uniforme convergentie}. (De
monotonie is onmisbaar: de glijdende bulten van
\cref{exo:b2:funcseq:2} convergeren puntsgewijs op een \omterm{def:b2:metric:compact}{compacte}
verzameling zonder uniformiteit.)
\end{solution}

\begin{solution}{exo:b2:funcseq:6}
Substitueer $u = nx$:
\[
\int_0^1 \frac{n f(x)}{1 + n^2x^2}\dd x
= \int_0^n \frac{f(u/n)}{1 + u^2}\,\dd u .
\]
De integranden $h_n(u) = \frac{f(u/n)}{1+u^2}\mathbf{1}_{u \leq n}$
convergeren puntsgewijs naar $\frac{f(0)}{1+u^2}$ (\omterm{def:b2:metric:continuity}{continuïteit} van
$f$ in $0$) en worden gedomineerd door $\frac{\norm f_\infty}{1 +
u^2}$, integreerbaar op $\intco{0}{\infty}$: de gedomineerde
convergentie (\cref{thm:b2:integration:dominated}) geeft als
limiet
\[
\int_0^\infty \frac{f(0)}{1 + u^2}\dd u = \frac{\pi}{2} f(0) .
\]
(De kernen concentreren zich bij $0$: een benaderende eenheid.)
\end{solution}

\begin{solution}{exo:b2:funcseq:7}
Gebruik voor $f(x) = x^2$ de tweede familie binomiale
identiteiten uit het bewijs: $\sum_k k^2 p_k = n(n-1)x^2 + nx$.
Bijgevolg is
\[
B_n(f)(x) = \sum_k \frac{k^2}{n^2}\,p_k
= \frac{n(n-1)x^2 + nx}{n^2}
= x^2 + \frac{x(1 - x)}{n} :
\]
dus $B_n(f) - f = \frac{x(1-x)}{n}$, met supnorm $\frac{1}{4n} =
O\bigl(\frac1n\bigr)$, zoals voorspeld.
\end{solution}

\begin{solution}{exo:b2:funcseq:8}
Voor $\varepsilon = 1$ is er een $N$ met $\norm{P_n - P_m}_{\infty,
\R} \leq 1$ voor $m, n \geq N$. Een veelterm die op $\R$ begrensd
is, is constant (een niet-constante gaat naar $\pm\infty$): dus
$P_n - P_m = c_{n,m}$, constanten. Voor $n \geq N$ is dus $P_n =
P_N + c_n$ met $c_n = P_n(0) - P_N(0)$ convergent (\omterm{def:b2:funcseq:def}{puntsgewijze
convergentie} in $0$). Bijgevolg is $f = \lim P_n = P_N + \lim c_n$:
een veelterm.
\end{solution}

\begin{solution}{exo:b2:funcseq:9}
(a) Er geldt $\norm{(3/4)^n\varphi(4^n\cdot)}_\infty = \frac12
(3/4)^n$: \omterm{def:b2:funcseq:series}{normale convergentie}, dus is $W$ \omterm{def:b2:metric:continuity}{continu}
(\cref{thm:b2:funcseq:seriestransfer}).

(b) Leg $x$ en $m$ vast; kies het teken van $h_m = \pm\frac12
4^{-m}$ zo dat het segment $\intcc{4^mx}{4^m(x + h_m)}$ (van lengte
$\frac12$) geen half geheel getal bevat, zodat $\varphi$ er affien
is met helling $\pm1$ (mogelijk: een interval van lengte $\frac12$
ontmoet hoogstens één half geheel getal; kies de kant die het
vermijdt).

Voor $n > m$: $4^n h_m = \pm\frac12 4^{n-m}$ is een geheel getal en
$\varphi$ is $1$-periodiek: de $n$-de term van het verschil
verdwijnt dus.

Voor $n = m$: $\abs{\varphi(4^m x + 4^m h_m) - \varphi(4^m x)} =
\abs{4^m h_m} = \frac12$ ($\varphi$ is op het segment affien met
helling $\pm 1$), dus draagt die term precies $(3/4)^m \cdot
\frac{1/2}{\abs{h_m}} = (3/4)^m\,4^m = 3^m$ bij aan het
quotiënt.

Voor $n < m$: de $1$-Lipschitz-eigenschap van $\varphi$ geeft
$\bigl|(3/4)^n\bigl(\varphi(4^nx + 4^nh_m) -
\varphi(4^nx)\bigr)\bigr| \leq (3/4)^n 4^n\abs{h_m} =
3^n\abs{h_m}$: elk draagt dus hoogstens $3^n$ bij aan het
quotiënt.

Bijgevolg is
\[
\Bigl|\frac{W(x + h_m) - W(x)}{h_m}\Bigr|
\geq 3^m - \sum_{n=0}^{m-1} 3^n
= 3^m - \frac{3^m - 1}{2} = \frac{3^m + 1}{2}
\longrightarrow \infty .
\]
Was $W$ in $x$ differentieerbaar, dan zou elk differentiequotiënt
langs $h_m \to 0$ naar $W'(x)$ convergeren: tegenspraak. $W$ is dus
overal \omterm{def:b2:metric:continuity}{continu} en nergens differentieerbaar.
\end{solution}

\begin{solution}{exo:b2:funcseq:10}
Puntsgewijs: voor $x \in \intco{0}{1}$ is de reeks meetkundig met
reden $-x$,
\[
\sum_{n\geq0}(-1)^n x^n(1-x) = \frac{1-x}{1+x},
\]
en in $x = 1$ verdwijnt elke term: som $0 = \frac{1-1}{2}$,
consistent --- de som is \omterm{def:b2:metric:continuity}{continu} op $\intcc{0}{1}$. Uniformiteit:
de restsom is een meetkundige staart,
\[
\abs{R_N(x)} = \frac{x^{N+1}(1-x)}{1+x} \leq x^{N+1}(1 - x)
\leq \max_{\intcc01} t^{N+1}(1-t)
= \frac{1}{N+2}\Bigl(\frac{N+1}{N+2}\Bigr)^{\!N+1}
\leq \frac{1}{N+2} \to 0 ,
\]
\omterm{def:b2:funcseq:def}{uniform} in $x$. Niet \omterm{def:b2:funcseq:series}{normaal}: $\norm{u_n}_\infty = \max x^n(1-x) =
\frac{1}{n+1}\bigl(\frac{n}{n+1}\bigr)^n \sim \frac{1}{\eu\,n}$, en
$\sum \frac1{\eu n}$ divergeert. Ter vergelijking: $\sum x^n(1-x)$
heeft partiële sommen $1 - x^{N+1}$, die puntsgewijs naar de
\emph{discontinue} $\mathbf 1_{\intco01}$ convergeren; volgens
\cref{thm:b2:funcseq:continuity} kan die convergentie op
$\intcc{0}{1}$ niet \omterm{def:b2:funcseq:def}{uniform} zijn.
\end{solution}

\begin{solution}{exo:b2:funcseq:11}
De limiet $f$ is \omterm{def:b2:metric:continuity}{continu} (\cref{thm:b2:funcseq:continuity}). Dan
is
\[
\abs{f_n(x_n) - f(x)}
\leq \abs{f_n(x_n) - f(x_n)} + \abs{f(x_n) - f(x)}
\leq \norm{f_n - f}_\infty + \abs{f(x_n) - f(x)} ,
\]
en beide termen gaan naar $0$ (\omterm{def:b2:funcseq:def}{uniforme convergentie}; \omterm{def:b2:metric:continuity}{continuïteit}
van $f$ in $x$). Tegenvoorbeeld bij louter \omterm{def:b2:funcseq:def}{puntsgewijze
convergentie}: $g_n(x) = nx\,\eu^{-nx^2} \to 0$ puntsgewijs op
$\intcc{0}{1}$ met $g_n$ en de limiet \omterm{def:b2:metric:continuity}{continu}, en toch is in $x_n =
\frac{1}{\sqrt{2n}} \to 0$
\[
g_n(x_n) = \sqrt{\frac n2}\,\eu^{-1/2} \longrightarrow +\infty
\neq 0 = f(0) .
\]
\end{solution}

\begin{solution}{exo:b2:funcseq:12}
\begin{enumerate}
  \item Inductie. Het geval $n = 1$ is de definitie. Neem de
        formule voor $n$ aan en zet $g(x) = \int_0^x
        \frac{(x-t)^n}{n!} f(t)\dd t$. Voor een integrand die
        \omterm{def:b2:metric:continuity}{continu} is in $(x,t)$ en $C^1$ in $x$, differentieert de
        \omterm{thm:b2:integration:continuity}{parameterintegraal} met variabele grens als
        \[
        g'(x) = \frac{(x-x)^n}{n!}f(x)
        + \int_0^x \frac{(x-t)^{n-1}}{(n-1)!}f(t)\dd t
        = T^nf(x)
        \]
        (splits $g(x+h) - g(x)$ in de strook $\int_x^{x+h}$, die
        $O(h\cdot\sup)$ is omdat de integrand in $t = x$ als $h^n$
        verdwijnt, en de vaste integraal van de toename in $x$, die
        met de middelwaardeongelijkheid en de \omterm{def:b2:metric:continuity}{continuïteit} wordt
        behandeld). Bovendien is $(T^{n+1}f)' = T^nf$
        (hoofdstelling van de integraalrekening) en $g(0) =
        T^{n+1}f(0) = 0$: twee primitieven van $T^nf$ die in $0$
        verdwijnen vallen samen, dus $T^{n+1}f = g$. De grens:
        \[
        \abs{T^nf(x)} \leq \norm f_\infty
        \int_0^x \frac{(x-t)^{n-1}}{(n-1)!}\dd t
        = \norm f_\infty\,\frac{x^n}{n!}
        \leq \frac{\norm f_\infty}{n!} .
        \]
  \item Er geldt $\sum_n \norm{T^nf}_\infty \leq \eu\,\norm
        f_\infty$: normale en dus \omterm{def:b2:funcseq:def}{uniforme convergentie}; $S$ is
        \omterm{def:b2:metric:continuity}{continu}. De partiële sommen voldoen aan $S_N = f + T
        S_{N-1}$, en $T$ is $1$-Lipschitz voor $\norm\cdot_\infty$
        (want $\abs{Tg(x)} \leq x\norm g_\infty$): met $N \to
        \infty$ in beide leden volgt $S = f + TS$.
  \item Zet $V(x) = f(x) + \eu^x\int_0^x \eu^{-t}f(t)\dd t$. Dan is
        $V - f$ van klasse $C^1$ met $(V-f)'(x) =
        \eu^x\int_0^x\eu^{-t}f + f(x) = V(x)$, en $(TV)' = V$ met
        $(V - f)(0) = TV(0) = 0$: dus $V - f = TV$, dat wil zeggen
        $V$ lost de vergelijking op. Eenduidigheid: zijn $S_1,
        S_2$ \omterm{def:b2:metric:continuity}{continue} oplossingen, dan voldoet $D = S_1 - S_2$ aan
        $D = TD$, dus $D = T^nD$ voor alle $n$ en $\norm D_\infty
        \leq \frac{\norm D_\infty}{n!} \to 0$: dus $D = 0$.
        Bijgevolg is
        \[
        \sum_{n\geq0} T^nf(x) = f(x) +
        \int_0^x \eu^{x-t}f(t)\,\dd t .
        \]
        (De reeks $\sum T^n$ is een meetkundige reeks van
        operatoren: een eerste voorproefje van de resolvente
        $(\mathrm{Id} - T)^{-1}$, uitgewerkt in het volume van
        bachelorjaar 3.)
\end{enumerate}
\end{solution}

\begin{solution}{pb:b2:funcseq:1}
\textbf{1.} De lineariteit is uit de formule duidelijk.
Positiviteit: de gewichten $p_k(x) \geq 0$, dus dwingt $f \geq 0$
af dat $B_nf \geq 0$; de monotonie volgt door dit op $g - f$ toe te
passen. Grens: uit $\pm f \leq \norm f_\infty$ volgt $\pm B_nf \leq
\norm f_\infty B_ne_0 = \norm f_\infty$. Eindpunten: $p_k(0) =
\mathbf 1_{k=0}$ en $p_k(1) = \mathbf 1_{k=n}$, dus $B_nf(0) =
f(0)$ en $B_nf(1) = f(1)$.

\textbf{2.} Differentieer $(x+y)^n = \sum_k\binom nk x^ky^{n-k}$
naar $x$, vermenigvuldig met $x$ en zet $y = 1 - x$:
\[
nx = \sum_k k\,p_k(x) ;
\]
twee keer, met vermenigvuldiging met $x^2$: $n(n-1)x^2 = \sum_k
k(k-1)p_k(x)$. Bijgevolg is $B_ne_0 = 1$ (binomium), $B_ne_1(x) =
\frac{nx}{n} = x$, en
\[
B_ne_2(x) = \frac{\sum_k k^2p_k}{n^2}
= \frac{n(n-1)x^2 + nx}{n^2}
= x^2 + \frac{x(1-x)}{n} .
\]

\textbf{3.} Uitwerken geeft
\[
\sum_k\Bigl(\frac kn - x\Bigr)^{\!2} p_k
= B_ne_2(x) - 2x\,B_ne_1(x) + x^2
= \frac{x(1-x)}{n} .
\]
Cauchy--Schwarz met de splitsing $\abs{k/n - x}\sqrt{p_k} \cdot
\sqrt{p_k}$:
\[
\sum_k\Bigl|\frac kn - x\Bigr| p_k
\leq \Bigl(\sum_k\Bigl(\frac kn - x\Bigr)^2
p_k\Bigr)^{\!1/2}
= \sqrt{\frac{x(1-x)}{n}} \leq \frac{1}{2\sqrt n},
\]
met $x(1-x) \leq \frac14$.

\textbf{4.} De gewichten $p_k(x)$ zijn niet-negatief met som $1$ en
zwaartepunt $\sum_k \frac kn p_k(x) = x$ (vraag 2). De eindige
ongelijkheid van Jensen voor de convexe $f$ (inductie vanaf de
definitie met twee punten, volume van bachelorjaar 1) geeft
\[
f(x) = f\Bigl(\sum_k \frac kn\,p_k\Bigr)
\leq \sum_k f\Bigl(\frac kn\Bigr)p_k = B_nf(x) .
\]

\textbf{5.} Op $\{k : \abs{k/n - x} > \delta\}$ is
$\bigl(\frac{k/n - x}{\delta}\bigr)^2 > 1$, dus
\[
\sum_{\abs{k/n-x}>\delta} p_k
\leq \frac{1}{\delta^2}\sum_k\Bigl(\frac kn -
x\Bigr)^2p_k
= \frac{x(1-x)}{n\delta^2} \leq \frac{1}{4n\delta^2} .
\]
Lezing: $p_k(x)$ is de verdeling van een steekproeffrequentie
$S_n/n$ van $n$ muntworpen met scheefheid $x$; haar
verwachtingswaarde is $x$ en haar variantie $\frac{x(1-x)}n \to 0$,
en de formule is de ongelijkheid van Chebyshev: de massa
concentreert zich bij $x$, zodat $f$ ertegen middelen in de limiet
$f(x)$ reproduceert (\cref{ch:b2:genfun} maakt de woordenschat
officieel).

\textbf{6.} Er geldt $\omega \leq 2M < \infty$; de monotonie is
duidelijk (supremum over een grotere verzameling). Heine: $f$
\omterm{def:b2:metric:continuity}{continu} op een \omterm{def:b2:metric:compact}{compacte} verzameling is uniform \omterm{def:b2:metric:continuity}{continu}, en dat zegt
precies dat $\omega(\delta) \to 0$ als $\delta \to 0^+$.
Subadditiviteit: geldt $\abs{s - t} \leq \delta_1 + \delta_2$, dan
voldoet het punt $u$ op het segment $\intcc st$ op afstand
$\min(\delta_1, \abs{s-t})$ van $s$ aan $\abs{s-u} \leq \delta_1$
en $\abs{u-t} \leq \delta_2$, en is $\abs{f(s)-f(t)} \leq
\abs{f(s)-f(u)} + \abs{f(u)-f(t)}$. Itereren geeft
$\omega(p\delta) \leq p\,\omega(\delta)$ voor $p \in \N^*$; voor
$\lambda > 0$, met $p = \lceil\lambda\rceil \leq 1 + \lambda$:
$\omega(\lambda\delta) \leq \omega(p\delta) \leq p\,\omega(\delta)
\leq (1+\lambda)\omega(\delta)$.

\textbf{7.} Omdat $\sum p_k = 1$:
\[
\abs{B_nf(x) - f(x)}
= \Bigl|\sum_k\bigl(f(k/n) - f(x)\bigr)p_k\Bigr|
\leq \sum_k\omega\bigl(\abs{k/n - x}\bigr)p_k .
\]
Voor elke $k$ geeft vraag 6 met $\lambda = \abs{k/n - x}/\delta$
dat $\omega(\abs{k/n-x}) \leq \bigl(1 +
\frac{\abs{k/n-x}}\delta\bigr)\omega(\delta)$; sommeren tegen de
$p_k$ levert de hoofdschatting.

\textbf{8.} Vul de grens van vraag 3 in:
\[
\abs{B_nf(x) - f(x)} \leq \Bigl(1 +
\frac{1}{2\delta\sqrt n}\Bigr)\omega(\delta),
\]
\omterm{def:b2:funcseq:def}{uniform} in $x$; met $\delta = n^{-1/2}$ is de haak $\frac32$: dus
$\norm{B_nf - f}_\infty \leq \frac32\omega(n^{-1/2}) \to 0$ volgens
vraag 6 (Heine). Dit bewijst
\cref{thm:b2:funcseq:weierstrass} opnieuw, met een snelheid.

\textbf{9.} $L$-Lipschitz betekent $\omega(\delta) \leq L\delta$:
snelheid $\frac{3L}{2\sqrt n}$. $\alpha$-höldercontinu betekent
$\omega(\delta) \leq C\delta^\alpha$: snelheid $\frac{3C}{2}
n^{-\alpha/2}$.

\textbf{10.} Met $k\binom{2m}k = 2m\binom{2m-1}{k-1}$:
\[
\sum_{k=m+1}^{2m}k\binom{2m}k
= 2m\sum_{j=m}^{2m-1}\binom{2m-1}{j}
= 2m\cdot 2^{2m-2},
\]
omdat $j \mapsto 2m-1-j$ de verzameling $\{m,\dots,2m-1\}$
bijectief op $\{0,\dots,m-1\}$ afbeeldt, zodat de som de helft van
$2^{2m-1}$ is. Ook is $\sum_{k=m+1}^{2m}\binom{2m}k = \frac{2^{2m}
- \binom{2m}m}{2}$ (dezelfde symmetrie). Bijgevolg
\[
\sum_{k=m+1}^{2m}(k-m)\binom{2m}k
= m\,2^{2m-1} - m\,\frac{2^{2m} - \binom{2m}m}{2}
= \frac m2\binom{2m}m .
\]
De substitutie $k \mapsto 2m-k$ beeldt de termen met $k < m$ af op
die met $k > m$ (gelijke binomiaalcoëfficiënten, gelijke
$\abs{k-m}$): de absolute som is dus twee keer de eenzijdige som,
namelijk $m\binom{2m}m$.

\textbf{11.} In $x = \frac12$ is $p_k(\tfrac12) =
\binom{2m}k2^{-2m}$ en $f(\tfrac12) = 0$, dus
\[
B_{2m}f\Bigl(\frac12\Bigr)
= \sum_k\Bigl|\frac{k}{2m} - \frac12\Bigr|
\binom{2m}k 2^{-2m}
= \frac{2^{-2m}}{2m}\,m\binom{2m}m
= \frac{\binom{2m}m}{2\cdot4^m}
\sim \frac{1}{2\sqrt{\pi m}}
\]
volgens \cref{ex:b2:comparison:centralbinomial}. Omdat hier
$\omega_f(\delta) = \delta$ is (de functie is $1$-Lipschitz en de
grens wordt aangenomen), voorspelt vraag 8 hoogstens
$\frac32(2m)^{-1/2}$: de werkelijke fout $\frac{1}{2\sqrt{\pi m}}$
heeft precies de orde $n^{-1/2}$ --- de snelheid is dus scherp op
de constante na.

\textbf{12.} Uit $f \leq g$ volgt $g - f \geq 0$, dus $L(g-f) \geq
0$, dat wil zeggen $Lf \leq Lg$. Uit $-\abs f \leq f \leq \abs f$
volgt $-L\abs f \leq Lf \leq L\abs f$, dat wil zeggen $\abs{Lf}
\leq L\abs f$.

\textbf{13.} Kies met Heine een $\delta$ zodanig dat
$\abs{f(s)-f(x)} \leq \varepsilon$ zodra $\abs{s-x} \leq \delta$.
Is $\abs{s - x} > \delta$, dan is $\frac{(s-x)^2}{\delta^2} > 1$ en
$\abs{f(s)-f(x)} \leq 2M \leq \frac{2M}{\delta^2}(s-x)^2$. In beide
gevallen geldt de beweerde grens.

\textbf{14.} Leg $x$ vast; vraag 13 zegt, als functies van $s$:
\[
-\varepsilon e_0 - \frac{2M}{\delta^2}q_x
\;\leq\; f - f(x)e_0
\;\leq\; \varepsilon e_0 + \frac{2M}{\delta^2}q_x,
\qquad q_x = e_2 - 2x\,e_1 + x^2e_0 .
\]
Pas de monotone lineaire $L_n$ toe (vraag 12) en evalueer in $x$:
\[
\abs{L_nf(x) - f(x)L_ne_0(x)}
\leq \varepsilon L_ne_0(x)
+ \frac{2M}{\delta^2}\bigl(L_ne_2(x) - 2xL_ne_1(x)
+ x^2L_ne_0(x)\bigr) .
\]

\textbf{15.} Schrijf $\alpha_j = L_ne_j - e_j$, zodat
$\norm{\alpha_j}_\infty \to 0$. Omdat $e_2(x) - 2xe_1(x) +
x^2e_0(x) = 0$, is
\[
L_ne_2(x) - 2xL_ne_1(x) + x^2L_ne_0(x)
= \alpha_2(x) - 2x\,\alpha_1(x) + x^2\alpha_0(x),
\]
met supnorm hoogstens $\norm{\alpha_2} + 2\norm{\alpha_1} +
\norm{\alpha_0} \to 0$. Ook geldt $L_ne_0 \to e_0$ \omterm{def:b2:funcseq:def}{uniform}, dus
$L_ne_0 \leq 2$ voor grote $n$, en $\abs{f(x)}\abs{L_ne_0(x) - 1}
\leq M\norm{\alpha_0} \to 0$. Samengevoegd met vraag 14 geeft dit
voor grote $n$, \omterm{def:b2:funcseq:def}{uniform} in $x$,
\[
\abs{L_nf(x) - f(x)} \leq 2\varepsilon +
\frac{2M}{\delta^2}\,o(1) + M\,o(1) \leq 3\varepsilon :
\]
dus $L_nf \to f$ \omterm{def:b2:funcseq:def}{uniform} --- de stelling van Korovkin.

\textbf{16.} Er geldt $B_ne_0 = e_0$ en $B_ne_1 = e_1$ exact, en
$\norm{B_ne_2 - e_2}_\infty = \max_x\frac{x(1-x)}{n} = \frac{1}{4n}
\to 0$ (vraag 2): Korovkin is dus van toepassing, en Weierstrass
volgt voor de derde keer.

\textbf{17.} $I_nf$ is lineair in $f$ (de knoopwaarden zijn dat),
en op elke cel is de affiene interpolant van niet-negatieve
knoopwaarden niet-negatief: dus positief. Verder is $I_ne_0 = e_0$
en $I_ne_1 = e_1$, want een affiene functie is haar eigen
interpolant. Op een cel $\intcc ab$ (met $b - a = \frac1n$) is de
affiene interpolant van $e_2$ gelijk aan $L(t) = (a+b)t - ab$, en
\[
L(t) - t^2 = (t-a)(b-t) \in
\intcc{0}{\tfrac{(b-a)^2}{4}} ,
\]
met het maximum in het midden: dus $\norm{I_ne_2 - e_2}_\infty =
\frac{1}{4n^2} \to 0$. Korovkin: $I_nf \to f$ \omterm{def:b2:funcseq:def}{uniform} voor elke
\omterm{def:b2:metric:continuity}{continue} $f$ --- veelhoekige benadering, zonder verdere schatting.

\textbf{18.} Bij gegeven $f$ en $\varepsilon$: Weierstrass levert
een veelterm $P = \sum_{j=0}^d a_jx^j$ met $\norm{f - P}_\infty
\leq \frac\varepsilon2$; elke $a_j$ vervangen door een rationale
$b_j$ met $\abs{a_j - b_j} \leq \frac{\varepsilon}{2(d+1)}$
verschuift de supnorm op $\intcc{0}{1}$ hoogstens
$\frac\varepsilon2$. De verzameling veeltermen met rationale
coëfficiënten is een \omterm{def:b2:structures:countable}{aftelbare} vereniging (over $d$) van \omterm{def:b2:structures:countable}{aftelbare
verzamelingen}, dus \omterm{def:b2:structures:countable}{aftelbaar}, en dicht: $C(\intcc{0}{1})$ is dus
separabel.

\textbf{19.} Wegens de lineariteit is $\int_0^1 fP = 0$ voor elke
veelterm $P$. Kies veeltermen $P_n \to f$ \omterm{def:b2:funcseq:def}{uniform} (Weierstrass):
\[
\Bigl|\int_0^1 f^2\Bigr|
= \Bigl|\int_0^1 f\,(f - P_n)\Bigr|
\leq \norm f_\infty\,\norm{f - P_n}_\infty
\longrightarrow 0 ,
\]
dus $\int_0^1 f^2 = 0$. Was $f(x_0) \neq 0$, dan geeft de
\omterm{def:b2:metric:continuity}{continuïteit} $f^2 \geq c > 0$ op een deelinterval, in strijd met
de verdwijnende integraal: dus $f = 0$. Bijgevolg vallen twee
\omterm{def:b2:metric:continuity}{continue} functies met dezelfde momenten $\int f t^n$ samen.

\textbf{20.} Met $p_{n,k}(x) = \binom nk x^k(1-x)^{n-k}$, de
afspraken $p_{n-1,-1} = p_{n-1,n} = 0$, de productregel en
$k\binom nk = n\binom{n-1}{k-1}$, $(n-k)\binom nk =
n\binom{n-1}{k}$ volgt
\[
p_{n,k}'(x) = n\bigl(p_{n-1,k-1}(x) - p_{n-1,k}(x)\bigr) .
\]
Sommeren tegen $f(k/n)$ en de index in de eerste som verschuiven
(Abelsommatie) geeft
\[
(B_nf)'(x) = n\sum_{j=0}^{n-1}\Bigl(f\Bigl(\frac{j+1}n\Bigr)
- f\Bigl(\frac jn\Bigr)\Bigr)p_{n-1,j}(x) .
\]

\textbf{21.} Volgens de middelwaardestelling is $f(\frac{j+1}n) -
f(\frac jn) = \frac1n f'(\xi_j)$ met $\xi_j \in
\intoo{j/n}{(j+1)/n}$, dus $(B_nf)'(x) = \sum_j
f'(\xi_j)\,p_{n-1,j}(x)$. Ook het knooppunt $\frac{j}{n-1}$ ligt in
$\intcc{j/n}{(j+1)/n}$ (beide ongelijkheden herleiden zich tot $j
\leq n-1$), dus $\abs{\xi_j - \frac j{n-1}} \leq \frac1n$ en
\[
\bigl|(B_nf)'(x) - B_{n-1}(f')(x)\bigr|
\leq \sum_j\Bigl|f'(\xi_j) -
f'\Bigl(\frac{j}{n-1}\Bigr)\Bigr| p_{n-1,j}(x)
\leq \omega_{f'}\Bigl(\frac1n\Bigr) \longrightarrow 0
\]
\omterm{def:b2:funcseq:def}{uniform}. Omdat $B_{n-1}(f') \to f'$ \omterm{def:b2:funcseq:def}{uniform}
(\cref{thm:b2:funcseq:weierstrass} toegepast op de \omterm{def:b2:metric:continuity}{continue} $f'$),
geeft de driehoeksongelijkheid $(B_nf)' \to f'$ \omterm{def:b2:funcseq:def}{uniform}. De
veeltermen $P_n = B_nf$ convergeren dan in $C^1$-zin naar $f$.

\textbf{22.} Taylor--Lagrange in $x$: $f(\frac kn) - f(x) =
f'(x)(\frac kn - x) + \frac{f''(\xi_k)}2(\frac kn - x)^2$.
Sommeren tegen $p_k$ doodt de lineaire term (vraag 2):
\[
\abs{B_nf(x) - f(x)}
\leq \frac{\norm{f''}_\infty}{2}\sum_k\Bigl(\frac kn -
x\Bigr)^2p_k
= \frac{\norm{f''}_\infty}{2}\cdot\frac{x(1-x)}{n}
\leq \frac{\norm{f''}_\infty}{8n} .
\]

\textbf{23.} Twee verdere differentiaties van $(x+y)^n$ geven de
factoriële momenten, met $n_{(j)} = n(n-1)\cdots(n-j+1)$:
\[
\sum_k k_{(j)}\,p_k = n_{(j)}\,x^j \qquad (j = 3, 4),
\]
en $k^3 = k_{(3)} + 3k_{(2)} + k$, $k^4 = k_{(4)} + 6k_{(3)} +
7k_{(2)} + k$ zetten die om in machtsmomenten:
\[
\sum_k k^3p_k = n_{(3)}x^3 + 3n_{(2)}x^2 + nx,
\qquad
\sum_k k^4p_k = n_{(4)}x^4 + 6n_{(3)}x^3 + 7n_{(2)}x^2 + nx .
\]
Werk $(k - nx)^4$ uit en verzamel (een geduldige maar zuiver
mechanische berekening met de vier machtsmomenten):
\[
\sum_k(k-nx)^4p_k = nx(1-x)\bigl(1 + 3(n-2)x(1-x)\bigr) .
\]
Met $x(1-x) \leq \frac14$ is het rechterlid hoogstens
$\frac n4\bigl(1 + \frac{3n}4\bigr) = \frac{3n^2}{16} + \frac n4
\leq n^2$ voor $n \geq 1$.

\textbf{24.} De vorm van Peano van Taylor in $x$ definieert
$\eta(t) = \frac{f(t) - f(x) - f'(x)(t-x) -
\frac12f''(x)(t-x)^2} {(t-x)^2}$ voor $t \neq x$ en $\eta(x) = 0$:
volgens Taylor--Lagrange is $\eta(t) = \frac12\bigl(f''(\xi) -
f''(x)\bigr)$ voor zekere $\xi$ tussen $t$ en $x$, dus $\abs\eta
\leq \norm{f''}_\infty$ en $\eta(t) \to 0$ als $t \to x$
(\omterm{def:b2:metric:continuity}{continuïteit} van $f''$). De ontwikkeling tegen $p_k$ sommeren en
de vragen 2--3 gebruiken geeft
\[
n\bigl(B_nf(x) - f(x)\bigr)
= \frac{x(1-x)}{2}f''(x)
+ n\sum_k\eta\Bigl(\frac kn\Bigr)\Bigl(\frac kn -
x\Bigr)^2p_k .
\]
Kies bij gegeven $\varepsilon$ een $\delta$ met $\abs\eta \leq
\varepsilon$ op $\abs{t - x}\leq\delta$. Nabij deel: hoogstens
$\varepsilon\,n\cdot\frac{x(1-x)}n \leq \varepsilon$. Ver deel: met
$C = \norm{f''}_\infty$ en vraag 23,
\[
n\,C\sum_{\abs{k/n-x}>\delta}\Bigl(\frac kn -
x\Bigr)^2p_k
\leq \frac{nC}{\delta^2}\sum_k\Bigl(\frac kn -
x\Bigr)^4p_k
= \frac{nC}{\delta^2 n^4}\sum_k(k-nx)^4p_k
\leq \frac{C}{\delta^2 n} \longrightarrow 0 .
\]
Bijgevolg gaat $n(B_nf(x) - f(x)) \to \frac{x(1-x)}2f''(x)$ --- de
stelling van Voronovskaja. Voor $f = e_2$ is dit bij elke $n$ exact
(\cref{exo:b2:funcseq:7}): het plafond $\frac1n$ is echt.

\textbf{25.} (i) De positiviteit zette \omterm{def:b2:funcseq:def}{puntsgewijze} ongelijkheden
om in operatorongelijkheden: zij leverde de normgrens, Jensen,
Chebyshev en heel Korovkin --- lineariteit alleen bewijst hier
niets. (ii) De \omterm{def:b2:metric:compact}{compactheid} kwam binnen via Heine (vragen 6 en 13),
via de begrensdheid van $f$, en doordat de \omterm{def:b2:nvs:norm}{norm}
$\norm\cdot_\infty$ überhaupt eindig is. (iii) Drie testfuncties
volstaan omdat de positiviteit alles herleidt tot het beheersen van
$L_n$ op de ene familie $(s-x)^2 = e_2 - 2xe_1 + x^2e_0$, waarvan
het opspansel dat van $e_0, e_1, e_2$ is. (iv) Bernstein
convergeert voor elke \omterm{def:b2:metric:continuity}{continue} $f$ met de eerlijke snelheid
$\omega(n^{-1/2})$ (scherp, vraag 11) maar verzadigt bij $\frac1n$
voor gladde $f$ (vraag 24); het fourierhoofdstuk voert hetzelfde
programma uit voor periodieke functies met de kern van Fejér ---
opnieuw een positieve operator met dezelfde deugden en dezelfde
bescheidenheid.
\end{solution}
